\documentclass[a4paper,11pt]{article}

\usepackage{amsmath, amssymb}
\usepackage{graphicx}
\usepackage{siunitx}
\usepackage[style=numeric, sorting=none]{biblatex}
\addbibresource{references.bib}

\usepackage{hyperref}
\hypersetup{
    colorlinks=true,
    linkcolor=blue,
    citecolor=blue,
    urlcolor=blue
}

\title{Neutrino Flavor Oscillations in the Presence of Primordial Black Holes: A Preprint Study}
\author{Autor: Your Name}
\date{\today}

\begin{document}

\maketitle

\begin{abstract}
This article investigates neutrino flavor oscillations influenced by the gravitational and matter fields of primordial black holes (PBHs). We develop an effective Hamiltonian framework incorporating local density profiles generated by PBHs of various scenarios, including stellar-origin and dark-matter related PBH populations. The formalism accounts for modifications in oscillation parameters due to these localized sources, with detailed derivation of the potential, effective mixing angles, and mass-squared differences. Numerical solutions to the evolution equations reveal characteristic resonance phenomena that depend on PBH distribution models. Our results offer new insights into neutrino signals that could probe early-universe PBH populations.
\end{abstract}

\section{Introduction}
Neutrino oscillations are sensitive probes of the fundamental properties of neutrinos and the matter through which they propagate. The presence of primordial black holes (PBHs) in the early universe could induce localized enhancements in matter density, modifying neutrino propagation. This study aims to model and analyze such effects, considering different scenarios of PBH distribution and mass spectrum.

\section{Methods}
\subsection{Effective Hamiltonian}
The flavor evolution of neutrinos in a medium affected by PBHs is described by the Schrödinger-like equation:
\begin{equation}
i \frac{d}{dr} \psi(r) = H(r, E) \psi(r),
\end{equation}
where the Hamiltonian
\begin{equation}
H(r, E) = H_\mathrm{vac}(E) + H_\mathrm{matt}(r),
\end{equation}
comprises vacuum oscillation terms and matter potentials influenced by PBH-induced density profiles.

\subsection{PBH Scenarios}
We consider three primary scenarios:
\begin{itemize}
    \item \textbf{Stellar-origin PBHs}: Mass range $10^{17}$ -- $10^{18}$ g, localized near density peaks.
    \item \textbf{Dark-sources PBHs}: Formation via early universe processes, distributed more uniformly, masses $10^{15}$ -- $10^{17}$ g.
    \item \textbf{Parameterized scenario}: General parametrization allowing variable density profiles from models.
\end{itemize}
The matter potential is modeled as:
\begin{equation}
V_\mathrm{LSC}(r) = \sqrt{2} G_F n_e(r),
\end{equation}
with $n_e(r)$ derived from PBH distribution functions.

\subsection{Numerical Solution}
The evolution equations are solved numerically using adaptive Runge-Kutta integrators. The resonance conditions are examined at different energies and radii, analyzing the behavior of the effective mixing parameters $\theta_\mathrm{eff}(r,E)$ and $\Delta m^2_\mathrm{eff}(r,E)$.

\section{Results}
Simulation results show that PBH populations induce localized resonances in the neutrino flavor conversion probability. The resonance positions depend on PBH mass spectrum and distribution, with potential observable signatures in neutrino telescopes.

\section{Discussion}
Our findings suggest that neutrino flavor data could serve as indirect probes of primordial black hole populations, especially in mass ranges affecting matter profile modulations significant enough to produce detectable resonance features.

\section{Conclusions}
The developed formalism provides a comprehensive framework for studying neutrino oscillations affected by PBHs. Future work will extend to include more detailed PBH clustering models and potential observational constraints.

\appendix

\section{Full Derivation of the Matter Potential}
The matter potential considering PBH-induced density fluctuations is given by:
\begin{equation}
V_\mathrm{LSC}(r) = \sqrt{2} G_F n_e^{\text{avg}}(r) + \delta V(r),
\end{equation}
where $\delta V(r)$ accounts for local density variations from PBHs.

\section{Sample Parameter Values}
\begin{table}[h]
\centering
\begin{tabular}{lcc}
\hline
Parameter & Default Value & Comment \\
\hline
$G_F$ & $1.1663787 \times 10^{-5}\,\mathrm{GeV}^{-2}$ & Fermi constant \\
$\Delta m^2$ & $2.5 \times 10^{-3}\,\mathrm{eV}^2$ & Mass-squared difference \\
$\theta_0$ & $33.82^\circ$ & Vacuum mixing angle \\
\hline
\end{tabular}
\caption{Default parameters used in simulations}
\end{table}

\bibliographystyle{unsrt}
\printbibliography

\end{document}
